% !TeX TS-program = xelatex

\documentclass[compress,aspectratio=169]{beamer}

\usepackage{presentationtemplate}
\usepackage[fontsize=\footnotesize, linenosfontsize=\tiny, askip=3mm]{mylisting}
\usepackage{myitemize}
\usepackage{emoji}

\title{Сравнение C и C++}

\begin{document}

\frame[plain]{\titlepage}

\begin{frame}

    \frametitle{Зачем C, если есть C++}

    \begin{myitemize}

        \item C проще, чем C++: нет классов, исключений,
            конструкторов и перегрузки.

        \item Рантайм C компактнее и универсальнее, чем
            рантайм C++.

        \item На C написано большое число проектов. Примеры:
            ядро Linux, SQLite, Nginx, Git.

    \end{myitemize}

\end{frame}

\begin{frame}[fragile]

    \frametitle{Один и тот же struct в двух языках}

    \begin{columns}[T]

        \begin{column}{0.48\textwidth}

            \myinputlisting[minted language=c]
                {presentations/02-c-vs-cpp/}
                {struct_c.c}

        \end{column}

        \begin{column}{0.48\textwidth}

            \myinputlisting[minted language=cpp]
                {presentations/02-c-vs-cpp/}
                {struct_cpp.cpp}

        \end{column}

    \end{columns}

\end{frame}

\begin{frame}[fragile]

    \frametitle{Ключевое слово struct — обязательно}

    \begin{myitemize}

        \item В C ключевое слово \texttt{struct} входит в
            тип переменной.

        \item Распространено создание псевдонима типа
            через \texttt{typedef}.

    \end{myitemize}

    \begin{columns}[T]

        \begin{column}{0.48\textwidth}

            \begin{mytitledinplacelisting}[minted language=c]{C}
struct Point { int x, y; };

// struct обязателен
struct Point p;

typedef struct Point Point;
// создали псевдоним Point
Point q;
            \end{mytitledinplacelisting}

        \end{column}

        \begin{column}{0.48\textwidth}

            \begin{mytitledinplacelisting}[minted language=c++]{C++}
struct Point { int x, y; };

// в C++ псевдоним
// создавать необязательно
Point p;
            \end{mytitledinplacelisting}

        \end{column}

    \end{columns}

\end{frame}

\begin{frame}[fragile]

    \frametitle{\texttt{void*} и приведения типов}

    В C \texttt{void*} неявно приводится к любому указателю
    на объект (и обратно).
    В C++ для этого требуется именованное преобразование
    \texttt{static\_cast}.

    \begin{columns}[T]

        \begin{column}{0.48\textwidth}

            \myinputlisting[minted language=c]
                {presentations/02-c-vs-cpp/}
                {malloc_c.c}

        \end{column}

        \begin{column}{0.48\textwidth}

            \myinputlisting[minted language=c++]
                {presentations/02-c-vs-cpp/}
                {malloc_cpp.cpp}

        \end{column}

    \end{columns}

\end{frame}

\begin{frame}[fragile]

    \frametitle{\texttt{NULL} и \texttt{nullptr}, ссылки}

    \begin{myitemize}

        \item Нулевой указатель: в C — макрос \texttt{NULL},
        в C++ — типизированный \texttt{nullptr} (C++11).

        \item Ссылок в C нет.

    \end{myitemize}

    \begin{columns}[T]

        \begin{column}{0.48\textwidth}

            \begin{mytitledinplacelisting}[minted language=c]{C}
int *p = NULL;

// ссылок нет
void swap(int *a, int *b);
            \end{mytitledinplacelisting}

        \end{column}

        \begin{column}{0.48\textwidth}

            \begin{mytitledinplacelisting}[minted language=c++]{C++}
int *p = nullptr;

// ссылки есть
void swap(int &a, int &b);
            \end{mytitledinplacelisting}

        \end{column}

    \end{columns}

\end{frame}

\begin{frame}[fragile]

    \frametitle{Область видимости констант перечислений}

    \begin{myitemize}

        \item В C константы перечислений живут в области, где объявленo
        перечисление. Конфликт имен приведет к ошибке компиляции.

        \item В C++ область видимости констант ограничена перечислением.

    \end{myitemize}

    \begin{columns}[T]

        \begin{column}{0.48\textwidth}

            \begin{mytitledinplacelisting}[minted language=c++]{C}
enum Color { RED, GREEN };

// error: ‘RED’ redeclared
// as different kind of symbol
int RED = 5;
            \end{mytitledinplacelisting}

        \end{column}

        \begin{column}{0.48\textwidth}

            \begin{mytitledinplacelisting}[minted language=c++]{C++}
enum Color { RED, GREEN };

// нет ошибки
int RED = 5;

enum class Color2 { RED };
int a = Color2::RED;
            \end{mytitledinplacelisting}

        \end{column}

    \end{columns}

\end{frame}

\begin{frame}[fragile]

    \frametitle{bool и псевдонимы типов}

    \begin{myitemize}

        \item В C типа \texttt{bool} нет в языке: только
        \texttt{\_Bool} + заголовок \texttt{<stdbool.h>} (C99),
        который вводит макросы \texttt{bool}, \texttt{true},
        \texttt{false}.

        \item В C++ \texttt{bool} — встроенный тип языка.

        \item Именованные типы: механизм один и тот же,
        в C пишут \texttt{typedef}, в C++ принято \texttt{using}.

    \end{myitemize}

    \begin{columns}[T]

        \begin{column}{0.48\textwidth}

            \begin{mytitledinplacelisting}[minted language=c]{C}
#include <stdbool.h>
// bool и true - макросы
// из stdbool.h
bool ok = true;

// псевдоним типа
typedef unsigned int uint;
            \end{mytitledinplacelisting}

        \end{column}

        \begin{column}{0.48\textwidth}

            \begin{mytitledinplacelisting}[minted language=c++]{C++}
// заголовочного файла
// не требуется
bool ok = true;

// почти равнозначно
using uint0 = unsigned int;
typedef unsigned int uint1;
            \end{mytitledinplacelisting}

        \end{column}

    \end{columns}

\end{frame}

\begin{frame}[fragile]

    \frametitle{Инициализация: значений по умолчанию в C нет}

    \begin{myitemize}

        \item В C поля структуры не имеют значений по умолчанию.

        \item В C++ для этого есть конструктор (начиная с C++11
            можно использовать default member initializers).

    \end{myitemize}

    \begin{columns}[T]

        \begin{column}{0.48\textwidth}

            \myinputlisting[minted language=c]
                {presentations/02-c-vs-cpp/}
                {config_c.c}

        \end{column}

        \begin{column}{0.48\textwidth}

            \myinputlisting[minted language=c++]
                {presentations/02-c-vs-cpp/}
                {config_cpp.cpp}

        \end{column}

    \end{columns}

\end{frame}

\begin{frame}[fragile]

    \frametitle{\texttt{malloc} и \texttt{new}}

    \begin{myitemize}

        \item \texttt{malloc} — функция: возвращает
            \texttt{void*}, при нехватке памяти — \texttt{NULL}.

        \item \texttt{new} — оператор: возвращает указатель
            нужного типа, при нехватке — исключение
            \texttt{std::bad\_alloc}.

        \item \texttt{malloc} возвращает указатель на выделенную
            память, \texttt{new} предварительно вызовет конструктор.

    \end{myitemize}

    \begin{columns}[T]

        \begin{column}{0.48\textwidth}

            \begin{mytitledinplacelisting}[minted language=c]{C++}
int *a = malloc(10 * sizeof *a);
if (!a) { /* обработка ошибки */ }
            \end{mytitledinplacelisting}

        \end{column}

        \begin{column}{0.48\textwidth}

            \begin{mytitledinplacelisting}[minted language=c++]{C}
// может бросить исключение
// std::bad_alloc
int *a = new int[10];
            \end{mytitledinplacelisting}

        \end{column}

    \end{columns}

\end{frame}

\begin{frame}[fragile]

    \frametitle{Перегрузка функций}

    Её в C нет \emoji{slightly-smiling-face}.

    \bigskip

    \begin{myinplacelisting}[minted language=c]
// ошибка компиляции
// error: conflicting types for ‘foo’
int foo();
void foo();
    \end{myinplacelisting}

\end{frame}

\begin{frame}[fragile]

    \frametitle{Префиксы вместо пространства имен}

    \begin{myitemize}

        \item В C конфликты имен решаются префиксами-конвенциями.

        \item В C++ — \texttt{namespace}: имена «упаковываются»
        в модуль, конфликты исключаются по построению.

    \end{myitemize}

    \begin{columns}[T]

        \begin{column}{0.48\textwidth}

            \begin{mytitledinplacelisting}[minted language=c]{C++}
int mylib_push(List *l, int v);
void mylib_free(List *l);
            \end{mytitledinplacelisting}

        \end{column}

        \begin{column}{0.48\textwidth}

            \begin{mytitledinplacelisting}[minted language=c++]{C}
namespace mylib {
bool push(List &l, int v);
void free(List &l);
}
            \end{mytitledinplacelisting}

        \end{column}

    \end{columns}

\end{frame}

\begin{frame}[fragile]

    \frametitle{Макросы для кодогонерации}

    \begin{myitemize}

        \item В C нет шаблонов. Кодогенерация возможна
            только засчет препроцессора.
            Макрос подставляет аргументы текстом.

        \item Шаблонная функция в C++ — обычный вызов: аргументы
            вычисляются один раз, тип выводится компилятором.

    \end{myitemize}

    \begin{columns}[T]

        \begin{column}{0.48\textwidth}

            \myinputlisting[minted language=c]
                {presentations/02-c-vs-cpp/}
                {max_c.c}

        \end{column}

        \begin{column}{0.48\textwidth}

            \myinputlisting[minted language=c++]
                {presentations/02-c-vs-cpp/}
                {max_cpp.cpp}

        \end{column}

    \end{columns}

\end{frame}

\begin{frame}[fragile]

    \frametitle{Ошибки: коды возврата и исключения}

    В C нет исключений. Функции могут сообщать об
    ошибке только через вовзрат значений или запись
    в определенный участок памяти.

    \begin{columns}[T]

        \begin{column}{0.48\textwidth}

            \myinputlisting[minted language=c]
                {presentations/02-c-vs-cpp/}
                {err_c.c}

        \end{column}

        \begin{column}{0.48\textwidth}

            \myinputlisting[minted language=c++]
                {presentations/02-c-vs-cpp/}
                {err_cpp.cpp}

        \end{column}

    \end{columns}

\end{frame}

\begin{frame}

    \frametitle{Что даёт знание C}

    \begin{myitemize}

        \item Контроль над памятью и понимание того, что происходит
            «под капотом» C++.

        \item Умение читать системный и легаси-код: ядра,
            драйверы, библиотеки.

        \item Понимание, чем C++ «платит» за удобство: каждая
            функциональность (конструкторы, исключения) имеет цену.

    \end{myitemize}

\end{frame}

\end{document}
